\subsection{An arithmetic successor preserving one original denominator}
\label{sec:exact-audit-fixed-y}

Let \(p=12h+1\) be prime, let
\(A_p=\{3h+1,\ldots,9h\}\), and choose \(a,a+1\in A_p\).
Retain the original residuals and supplies
\[
 b=a+1,\qquad R=4a-p,\qquad R_b=R+4,\qquad
 S_a=pa,\qquad S_b=p(a+1),\qquad A=a(a+1).
\]
For \(c\in\{a,b\}\), let
\[
 \mathcal I_c=\left\{(y,z)\in\mathbb Z_{>0}^2:
       \frac4p=\frac1c+\frac1y+\frac1z\right\}.
\]
Every retained coordinate, including the ordered unmarked denominators,
is part of these objects.

\begin{theorem}[Exact integral fixed-denominator successor]
\label{thm:exact-audit-fixed-y}
Define the two finite domains
\begin{align*}
 \mathcal D_a&=\{(y,z)\in\mathcal I_a:A+z\mid A^2\},\\
 \mathcal D_b^-&=\{(y,z')\in\mathcal I_b:
                       0<z'<A,\ A-z'\mid A^2\}.
\end{align*}
Then
\begin{equation}\label{eq:exact-audit-fixed-y-map}
 T_a:\mathcal D_a\longrightarrow\mathcal D_b^-,\qquad
 (y,z)\longmapsto\left(y,\frac{Az}{A+z}\right)
\end{equation}
is a bijection, with inverse
\begin{equation}\label{eq:exact-audit-fixed-y-inverse}
 (y,z')\longmapsto\left(y,\frac{Az'}{A-z'}\right).
\end{equation}
Every forward value satisfies \(0<z'<\min(z,A)\). Thus this
successor strictly decreases the original denominator \(z\) on its
complete integral-return domain; every fibre is a singleton.

For the source reduced positive primitive ratio \(\tau=m/n=y/z\),
the target ratio is
\begin{equation}\label{eq:exact-audit-fixed-y-ratio}
 \tau'=\frac{a(R+4)\tau+p}{(a+1)R}.
\end{equation}
Its raw integer pair, reduction factor, and primitive return are exactly
\begin{equation}\label{eq:exact-audit-fixed-y-gcd}
 U=a(R+4)m+pn,\quad V=(a+1)Rn,\quad
 g=\gcd(U,V),\quad (m',n')=(U/g,V/g).
\end{equation}
For every source witness, this primitive target obeys
\[
 m',n'\mid p(a+1),\qquad R+4\mid m'+n'
\]
if and only if the proved divisibility \(A+z\mid A^2\) holds.
In particular the reduction factor and all target exponent bounds are
retained, not replaced by projective-ratio existence.
\end{theorem}
\begin{proof}
Keeping the same original \(y\), subtraction of the equations for
shells \(a\) and \(b=a+1\) forces
\[
 \frac1{z'}=\frac1z+\frac1a-\frac1{a+1}
           =\frac1z+\frac1A.
\]
Its unique positive rational solution is \(z'=Az/(A+z)\).
Conversely this identity proves the original equation at \(b\),
because
\[
 \frac1b+\frac1y+\frac1{z'}
 =\frac1b+\frac1y+\frac1z+\frac1a-\frac1b=\frac4p.
\]
Since \(A,z>0\), both \(z'/z=A/(A+z)<1\) and
\(z'/A=z/(A+z)<1\), proving strict positivity and both bounds.
The exact identity
\begin{equation}\label{eq:exact-audit-fixed-y-factor}
 z'=A-\frac{A^2}{A+z},\qquad (A+z)(A-z')=A^2
\end{equation}
shows, for integer source \(z\), that \(z'\) is an integer exactly
when \(A+z\mid A^2\). In that case \(A-z'>0\) is a positive
integer divisor of \(A^2\); thus the forward point lies in
\(\mathcal D_b^-\).

Conversely, for a point in \(\mathcal D_b^-\), put
\[
 z=\frac{Az'}{A-z'}=\frac{A^2}{A-z'}-A.
\]
It is a positive integer by \(0<z'<A\) and the displayed divisor
condition. Its reciprocal satisfies \(1/z=1/z'-1/A\), so
\[
 \frac1a+\frac1y+\frac1z
 =\frac1a+\frac1y+\frac1{z'}-\frac1a+\frac1b=\frac4p.
\]
Also \(A+z=A^2/(A-z')\) divides \(A^2\). Thus it lies in
\(\mathcal D_a\). Equation \eqref{eq:exact-audit-fixed-y-factor}
proves both inverse compositions. The domains are finite: on the
source, \(A+z\) is a positive divisor of the fixed positive integer
\(A^2\), and each \(z\) fixes \(y\) uniquely through the original
equation; the target has \(1\le z'<A\), again with at most one
\(y\) for each \(z'\).

The source shell identity gives
\(y=pa(1+\tau)/R\), since \(y/z=\tau\) and
\(R/pa=1/y+1/z=(1+\tau)/y\). Consequently
\begin{align*}
 \tau'=\frac y{z'}
 &=\frac yz+\frac yA
 =\tau+\frac{p(1+\tau)}{(a+1)R}\\
 &=\frac{((a+1)R+p)\tau+p}{(a+1)R}
 =\frac{a(R+4)\tau+p}{(a+1)R},
\end{align*}
where the last equality uses the original identity \(R+p=4a\).
This proves \eqref{eq:exact-audit-fixed-y-ratio}; substitution
\(\tau=m/n\) and division by the explicitly retained \(g\) prove
\eqref{eq:exact-audit-fixed-y-gcd}. All four integers \(U,V,m',n'\)
are positive, and \(\gcd(m',n')=1\).

To verify the target budget claim without assuming any lift, recall its
elementary exact reconstruction. For an integer target witness,
\[
 d_b=(R+4)y-pb>0,\qquad e_b=(R+4)z'-pb>0,\qquad
 d_b e_b=(pb)^2.
\]
Positivity follows from \((R+4)/pb=1/y+1/z'\); the product follows
by expansion of the same identity. Moreover
\(d_b/(pb)=y/z'=m'/n'\). If \(v_\ell(pb)=e_\ell\), then
\(0\le v_\ell(d_b)\le2e_\ell\). Reducing \(d_b/(pb)\) gives
\[
 v_\ell(m')=\max(v_\ell(d_b)-e_\ell,0),\quad
 v_\ell(n')=\max(e_\ell-v_\ell(d_b),0),
\]
so both are at most \(e_\ell\). Thus \(m',n'\mid pb\).
Since \(\gcd(pb,R+4)=1\), reducing
\(d_b\equiv-pb\pmod{R+4}\) gives \(R+4\mid m'+n'\).

In the reverse direction those bounded primitive conditions reconstruct
positive integers
\[
 Y=\frac{pb}{n'}\frac{m'+n'}{R+4},\qquad
 Z=\frac{pb}{m'}\frac{m'+n'}{R+4}.
\]
Their ratio is \(m'/n'\), and their reciprocal sum is
\((R+4)/(pb)\), exactly the target rational shell. A positive
ratio \(m'/n'\) and that reciprocal sum determine the ordered rational
pair uniquely by these same formulas; hence \((Y,Z)=(y,z')\).
Thus \(z'\) is integral, which by
\eqref{eq:exact-audit-fixed-y-factor} is precisely \(A+z\mid A^2\).
This proves both directions of the budget equivalence without an
unproved arithmetic-return premise.
\end{proof}

At \(p=13,a=4,R=3\), the actual source pair \((m,n)=(13,2)\)
gives
\[
 (a,y,z)=(4,130,20),\qquad A=20,\qquad
 (A+z,A-z')=(40,10),\qquad (b,y,z')=(5,130,10).
\]
The raw target pair is
\((U,V)=(4\cdot7\cdot13+13\cdot2,\,5\cdot3\cdot2)
=(390,30)\); its \(g=30\) gives \((m',n')=(13,1)\).
The original identities can also be checked over denominators
\(260\) and \(130\): their reciprocal numerators are respectively
\(65+2+13=80\) and \(26+1+13=40\).

At \(p=37,a=12,R=11\), the actual full source witness
\((m,n)=(3,74)\) gives \((a,y,z)=(12,42,1036)\).
Here
\[
 A=156,\quad A+z=1192,\quad A^2=24336
 =20\cdot1192+496.
\]
Thus the exact positive rational successor has
\[
 z'=\frac{20202}{149},\quad
 (U,V)=(3278,10582),\quad g=22,\quad(m',n')=(149,481).
\]
Its target supply is \(p(a+1)=481=37\cdot13\), and the retained
numerator \(149\) does not divide it. Indeed \(481=3\cdot149+34\).
Even though \(15\mid149+481=630\), the actual divisor budget fails.
This computes the failure of this specific successor on an occupied
source shell, while the successful \(p=13\) case proves that its
integer-return domain is nonempty.
